% !TEX program = xelatex
% =====================================================================
%  高等数学习题课讲义 — LaTeX 模板
%  风格还原:北京大学高等数学B习题课讲义(楷体正文 / 黑体标题 / 定理环境)
%  格式参考来源:https://darkoxie.github.io/ (原讲义著作权归原作者所有)
%  编译:  latexmk -xelatex <file>.tex      (推荐,自动跑足遍数)
%         或  xelatex <file>.tex  跑两遍
% =====================================================================
\documentclass[11pt]{ctexart}

% ---------- 页面 ----------
\usepackage{geometry}
\geometry{a4paper, top=2.5cm, bottom=2.5cm, left=2.4cm, right=2.4cm}

% ---------- 中文字体:正文楷体,标题黑体 ----------
\setCJKmainfont{KaiTi}                 % 正文 → 楷体(还原参考讲义)
\ctexset{
  section/format       = \Large\bfseries\heiti,
  subsection/format    = \large\bfseries\heiti,
  subsubsection/format = \normalsize\bfseries\heiti,
}

% ---------- 数学 ----------
\usepackage{amsmath, amssymb, amsthm, mathtools}
\usepackage{bm}

% ---------- 图形 ----------
\usepackage{graphicx}
\usepackage{tikz}
\usepackage{pgfplots}
\pgfplotsset{compat=1.18}
\usetikzlibrary{calc, arrows.meta, intersections}

% ---------- 列表 / 超链接 ----------
\usepackage{enumitem}
\usepackage[hidelinks]{hyperref}

% ---------- 定理类环境(中文)----------
\theoremstyle{plain}
\newtheorem{theorem}{定理}[section]
\newtheorem{proposition}{命题}[section]
\newtheorem{lemma}{引理}[section]
\newtheorem{corollary}{推论}[section]

\theoremstyle{definition}
\newtheorem{definition}{定义}[section]
\newtheorem{example}{例}            % 连续编号:例 1, 例 2, ...
\newtheorem{exercise}{习题}

\theoremstyle{remark}
\newtheorem*{remark}{注}

\renewcommand{\proofname}{证明}      % proof 环境引导词 → “证明”,结尾自带 □

% ---------- 常用数学宏 ----------
\newcommand{\R}{\mathbb{R}}
\newcommand{\N}{\mathbb{N}}
\newcommand{\Z}{\mathbb{Z}}
\newcommand{\Q}{\mathbb{Q}}
\newcommand{\dd}{\,\mathrm{d}}
\newcommand{\eps}{\varepsilon}
\DeclareMathOperator{\sgn}{sgn}

% ---------- 标题信息(按需修改)----------
\title{\heiti 北京大学高等数学B习题课讲义:预备知识}
\author{作者姓名 \\ 北京大学数学科学学院}
\date{最后修改:\today}

% =====================================================================
\begin{document}
\maketitle

% ------- 引言段(可选)-------
本讲义大致按照教材顺序梳理本章的核心概念、定理与典型例题,供初学者复习与查阅之用。

% =========================== 正文从此开始 ===========================

\section{知识点详解}

\subsection{实数和不等式}

% 无编号的小标题(参考讲义中“有理数和无理数”这类)用带星号的 subsubsection
\subsubsection*{有理数和无理数}

有理数是可以写成两个整数之比 $\frac{p}{q}$($q\neq 0$)的数,全体有理数记作 $\Q$;
无法写成这种形式的实数称为无理数。

\begin{theorem}[算术—几何平均不等式]
设 $a_1,a_2,\dots,a_n\ge 0$,则
\begin{equation}
  \frac{a_1+a_2+\cdots+a_n}{n}\ \ge\ \sqrt[n]{a_1 a_2\cdots a_n},
\end{equation}
当且仅当 $a_1=a_2=\cdots=a_n$ 时取等号。
\end{theorem}

\begin{example}
求证 $\sqrt{2}$ 是无理数。
\end{example}

\begin{proof}
反证法。设 $\sqrt{2}=\frac{p}{q}$,其中 $p,q$ 互质。则 $p^2=2q^2$,故 $p$ 为偶数,
记 $p=2k$,代入得 $q^2=2k^2$,从而 $q$ 也为偶数,与 $p,q$ 互质矛盾。
\end{proof}

\begin{remark}
上述反证法的关键在于“互质”假设导出的奇偶性矛盾。
\end{remark}

% =========================== 正文到此结束 ===========================

\end{document}
